\section{Problem and Fusion Rule}
\subsection{Robot, sensing, and communication model}
There are $N$ robots with known poses $s_i(t)$ and a common clock. A physical
edge exists when two robots are within radio range. A geometry-only
coordinator selects a bidirectional minimum spanning tree in each connected
component. At each time, a robot predicts its LMB, updates with its own
measurement set, and exchanges the resulting posterior for one synchronous
fusion round. A dropped neighbor's Metropolis weight returns to the receiver's
self weight. Fusion output is retained as the prior for the next local step.
The coordinator is modeled as reliable and its traffic is charged; the
estimator messages are lossy. This is not a decentralized tree-agreement
protocol or a consensus-convergence guarantee.

A label $\ell=(t_b,b)$ specifies a shared birth time and candidate region.
For represented label $\ell$, source $j$ supplies
$\pi_j^\ell=(r_j^\ell,p_j^\ell)$, where $r_j^\ell\in[0,1]$ and
$p_j^\ell(x)$ is a Gaussian-mixture density over position and velocity.
The birth prior does not disclose whether a target actually exists.
All methods receive the same region/time prior and already-consistent keys.
They do not use truth cardinality, target trajectories, or evaluation
assignments in fusion decisions.

\subsection{Observation history and participation}
Each retained Bernoulli carries two metadata values. The lineage flag
$h_j^\ell$ records whether the density descends from any scheduled local
update with nonzero detection opportunity. A direct opportunity is evaluated
on the predicted Gaussian means before the measurement update:
\begin{equation}
 o_j^\ell(t)=\mathbf{1}\!\left[
   \sum_g w_{jg}^\ell p_{D,j}(\mu_{jg}^\ell,t)>0\right].
 \label{eq:opportunity}
\end{equation}
The local timestamp $\tau_j^\ell$ is the last time this condition held at
robot $j$. Both detections and missed detections count. A received posterior
can propagate lineage but cannot advance $\tau_j^\ell$: after fusion the
receiver attaches its own local timestamp. Consequently the age measures
\emph{direct sensor opportunity}, not the age of every datum in a correlated
posterior. We deliberately do not infer measurement novelty from packet
arrival or attempt to remove all common information.

The participation rule is a controlled baseline component. If at least one
input carries observation lineage, represented inputs without lineage are
excluded as untouched priors. Otherwise they retain ordinary participation.
An absent label contributes a negative existence censor only when the
represented spatial density places at least half its Gaussian-mean mass
inside that source's current FoV, with expected detection at least $0.1$.
Its existence upper bound is $0.01$. Other absent sources abstain. Such
censors are qualified negative evidence and are never assigned a spatial
density. The flags propagate only through actual participating inputs.
This rule uses a mean-based FoV approximation, not exact integration through
a sensor boundary.

Let $\mathcal P$ denote eligible represented inputs and $\mathcal E$ eligible
existence inputs, including qualified absence censors. For the receiver's
delivered Metropolis weights $w_j$, define
\begin{equation}
 a_j=\frac{w_j\mathbf1[j\in\mathcal P]}{\sum_{k\in\mathcal P}w_k},
 \qquad
 q_j=\frac{w_j f_j\mathbf1[j\in\mathcal E]}{\sum_{k\in\mathcal E}w_k f_k}.
 \label{eq:weights}
\end{equation}
For a represented label with a positive timestamp set
\begin{equation}
 f_j=\rho+(1-\rho)\exp\!\left[-(t-\tau_j^\ell)\Delta t/\lambda\right].
\end{equation}
Use $f_j=\rho$ for a represented but never-directly-observed label, and
$f_j=1$ for a qualified absence censor.
Here $\rho=0.25$ and $\lambda=5$ s are fixed before development outcomes.
The age factor is independent of the sign of the existence claim: a recent
low $r_j$ can outweigh an older high $r_j$. Positive self weight and the
existing absence logic supply the ordinary single-input fallback; a label
without eligible spatial support is omitted.

\subsection{Existence recency with fixed spatial weights}
For clarity suppress the label and receiver indices. Denote Bernoulli
relative entropy by $D_B$ and ordinary density relative entropy by $D$.
For fixed eligibility and normalized $a,q$, consider
\begin{equation}
 \min_{r\in(0,1),\,p}
 J(r,p)=\sum_{j\in\mathcal E}q_j D_B(r\|r_j)
       +r\sum_{j\in\mathcal P}a_j D(p\|p_j).
 \label{eq:objective}
\end{equation}
When the two supports and weight vectors agree, this is the Bernoulli KLA
objective. Allowing separate vectors changes that objective; ordinary KLA
consistency claims cannot simply be transferred.

\noindent\textbf{Proposition 1 (fixed-input form).}
Assume $0<r_j<1$, $\eta_a=\int\prod_{j\in\mathcal P}p_j(x)^{a_j}dx>0$,
and finite relevant divergences. A minimizer of~\eqref{eq:objective} is
\begin{align}
 p^*(x)&=\eta_a^{-1}\prod_{j\in\mathcal P}p_j(x)^{a_j},\label{eq:spatial}\\
 \operatorname{logit}r^*&=\sum_{j\in\mathcal E}q_j\operatorname{logit}r_j
                           +\log\eta_a.\label{eq:existence}
\end{align}
For fixed $a$, eligibility, and input densities, changing $q$ changes
existence without changing $p^*$.

\noindent\emph{Proof.}
Writing $p_a=\eta_a^{-1}\prod_jp_j^{a_j}$ gives
$\sum_j a_jD(p\|p_j)=D(p\|p_a)-\log\eta_a$.
For every $r>0$ this is minimized at $p=p_a$. Differentiating the remaining
scalar objective yields
$\operatorname{logit}r-\sum_jq_j\operatorname{logit}r_j-\log\eta_a=0$.
Its second derivative is $1/[r(1-r)]>0$, establishing the stated solution.
\hfill$\square$

The overlap term remains a penalty: for exact normalized densities,
$0<\eta_a\leq1$. Thus the rule does not guarantee cardinality consistency,
probability calibration, or an existence value between all inputs. Nor does
the proposition guarantee identical recursive localization: different
existence decisions can affect subsequent associations, pruning, and which
states are extracted. We test localization on common assignment support
in addition to global set metrics.

The implementation calls the existing mixture-aware Gaussian-power fusion
routine with $a$ for spatial pooling and $q$ for existence. It retains at
most three components per input for the powered-mixture calculation, at
most 256 component tuples, and at most eight output components. A
non-admissible approximate overlap normalizer triggers the existing
moment-matched fallback. Fractional powers
of Gaussian mixtures and component reduction are approximate, so the
proposition describes the ideal densities rather than an exact mixture
implementation. All arms retain the same eight-component cap, prediction,
measurement-update code, and MAP-cardinality extraction. We call the rule
\emph{existence recency} (ER); using $q$ in both blocks is the distinct
\emph{both-block recency} ablation.

\subsection{Confirmation ablation and communication cost}
A separate ablation delays lineage restriction until some input carries a
positive confirmation. A robot creates this bit after two consecutive local
updates with opportunity and detection-association mass at least $0.5$.
Empty measurements reset the local streak, even if the update retains old
diagnostic fields. Only the bit propagates; relays cannot supply consecutive
local hits or refresh direct age. Ordinary FoV-aware pooling applies before
confirmation. This conventional two-hit gate tests a false-positive versus
short-contact discovery tradeoff; it is not presumed beneficial.

All communicating arms serialize the full retained mixture and identical
metadata, then decode the packet before receiver fusion. Each attempted
packet occupies 16,384 bytes, with a fixed 128-byte-per-robot control charge
per frame. Padding isolates fusion at equal attempted wire cost. We also
report unpadded payloads because differing mixture sizes can change actual
data volume. The method adds constant metadata per represented label and
linear-time weight calculations; it does not reduce the scheduled traffic.
